\chapter{PF2 Ratio Majorants and a Larger Zero-Free Disk}
\label{chap:pf2-ratio-zero-free-disk}

\status{SOH-G019 proved from SOH-G005 PF2 and the entirety of the even quotient $F$: the closed disk $|w|\le a_0/(2a_1)$ is zero-free.}

\section{Purpose}

Chapter~\ref{chap:quarter-disk-zero-exclusion} established an explicit rational zero-free disk $|w|\le1/4$ and used it to close the negative-inversion paired-zero frontier. The present chapter uses the stronger structural information already available from PF2 to obtain a much larger zero-free disk.

Write
\[
 F(w)=\sum_{k\ge0}a_k w^k,
 \qquad a_k>0,
\]
with
\[
 \xi\!\left(\frac12+z\right)=F(z^2).
\]
SOH-G005 proves
\[
 a_k^2\ge a_{k-1}a_{k+1}
 \qquad(k\ge1).
\]

\section{Monotone coefficient ratios}

Define
\[
 q_k:=\frac{a_{k+1}}{a_k}.
\]
Since the coefficients are positive, PF2 is equivalent to
\[
 q_k\le q_{k-1}.
\]
Thus with
\[
 q_0=\frac{a_1}{a_0}
\]
we have
\[
 q_k\le q_0
 \qquad(k\ge0).
\]
Inductively,
\[
 a_k
 =a_0\prod_{j=0}^{k-1}q_j
 \le a_0q_0^k.
\]

\section{Geometric tail majorant}

For $r\ge0$ with $q_0r<1$,
\[
 \sum_{k\ge1}a_kr^k
 \le a_0\sum_{k\ge1}(q_0r)^k
 =a_0\frac{q_0r}{1-q_0r}.
\]
Define
\[
 \boxed{
 R_0:=\frac1{2q_0}=\frac{a_0}{2a_1}.}
\]
If $r<R_0$, then $q_0r<1/2$, and therefore
\[
 \sum_{k\ge1}a_kr^k<a_0.
\]
Hence
\[
 |F(w)|
 \ge a_0-\sum_{k\ge1}a_k|w|^k
 >0
 \qquad(|w|<R_0).
\]

\section{Strictness on the boundary}

At $r=R_0$ the geometric estimate gives
\[
 \sum_{k\ge1}a_kR_0^k\le a_0.
\]
Equality would require
\[
 a_k=a_0q_0^k
 \qquad\text{for every }k\ge1,
\]
because every term of the positive sum is bounded termwise by the corresponding geometric majorant.

But then
\[
 \limsup_{k\to\infty}a_k^{1/k}=q_0>0,
\]
so the power series of $F$ would have finite radius of convergence $1/q_0$. This contradicts the fact that $F$ is entire. Therefore at least one coefficient inequality is strict and
\[
 \sum_{k\ge1}a_kR_0^k<a_0.
\]
Consequently,
\[
 \boxed{
 F(w)\ne0\quad\text{for every }|w|\le R_0.}
\]

\section{Canonical numerical scale}

The first two coefficients are
\[
 a_0=F(0)=\xi(1/2)
\]
and
\[
 a_1=\frac12\frac{d^2}{dz^2}\xi(1/2+z)\bigg|_{z=0}.
\]
High-precision regression gives
\[
 a_0\approx0.4971207781883141,
 \qquad
 a_1\approx0.01148597215757272,
\]
so
\[
 \boxed{R_0\approx21.64034403742489.}
\]
The numerical value is diagnostic only; the theorem is symbolic.

\section{Consequence in the critical-axis coordinate}

Because
\[
 w=z^2,
 \qquad z=s-\frac12,
\]
every zero $\rho$ of $\xi$ satisfies
\[
 \left|\rho-\frac12\right|^2>R_0.
\]
Thus
\[
 \boxed{
 \left|\rho-\frac12\right|>\sqrt{R_0}.}
\]
Numerically,
\[
 \sqrt{R_0}\approx4.6519.
\]
This gives a zero-free Euclidean disk around the half-axis origin in the centered $s$-plane.

\section{Relation to the preceding no-go chain}

G018 supplies a completely explicit rational disk and directly proves that negative inversion cannot map any xi zero to another xi zero. G019 is a stronger downstream consequence of PF2: it replaces the radius $1/4$ by the symbolic coefficient scale $a_0/(2a_1)$.

The two results therefore serve different evidential roles. G018 is elementary and explicit; G019 is structural and stronger.

\section{Proof firewall}

G019 proves coefficient-ratio monotonicity from PF2, the geometric coefficient majorant, the closed zero-free disk $|w|\le a_0/(2a_1)$, and the corresponding centered $s$-plane zero-free disk. It does not prove real-rootedness of $F$, PF$\infty$, SOH-G003 real-rootedness, or the Riemann Hypothesis.
